% ============================================================
%  Chapter 7 — Conclusion
% ============================================================
\chapter{Conclusion}
\label{chap:conclusion}

\section{Summary of Findings}

This thesis investigated whether weather variables significantly improve day-ahead electricity price forecasting for the French market, and whether the resulting forecasts generate economic value through a directional trading strategy. Our analysis, conducted on a comprehensive hourly dataset spanning January 2018 to April 2025, leads to four principal conclusions.

\textbf{First, machine learning models substantially outperform the na\"{i}ve benchmark.} Both Random Forest variants reduce MAE by approximately 49\% (from 33.09 to $\approx$16.94 EUR/MWh) and increase R$^2$ from 0.16 to 0.78 on the out-of-sample test set (May 2024--April 2025). The Diebold-Mariano test confirms this improvement is highly statistically significant ($p < 0.001$). The performance is robust: both RF models record positive monthly Sharpe ratios in every month of the test period.

\textbf{Second, the role of weather is regime-dependent.} In the stable
May 2024--April 2025 test period, weather features produce no statistically
significant forecast improvement ($p = 0.572$, DM $= -0.565$). However, a
validation on the 2022 energy crisis reverses this result: the DM test yields
$p < 0.001$ (DM $= -13.27$), and Model C reduces MAE by 0.73 EUR/MWh over
Model B in 2022. We attribute the stable-period non-significance to the
\textit{information redundancy hypothesis}: the load forecast, TTF gas price,
and nuclear availability ratio already encode the meteorological information
relevant to French price formation under normal conditions. The significance
in 2022 arises because the gas supply crisis decoupled TTF from local
temperature, breaking the indirect weather channel and making the direct
temperature-demand effect visible. This regime-dependent finding is the thesis's
primary empirical contribution and calls for nuanced interpretation: weather
features are neither always beneficial nor always redundant — their value
depends on the structural state of the gas and nuclear supply.

\textbf{Third, Random Forest outperforms XGBoost with default hyperparameters.} XGBoost achieves a higher absolute P\&L in the backtest but at the cost of a Maximum Drawdown more than three times larger than the Random Forest. On a risk-adjusted basis (Calmar ratio of 107 vs 320--328), RF is the superior model. The performance gap is attributed to XGBoost's sensitivity to extreme price observations in the training set and could be closed through systematic hyperparameter optimisation.

\textbf{Fourth, the forecasting models generate economically significant value.} Under the central transaction cost scenario (0.30 EUR/MWh), RF models generate net P\&L of approximately 136,700 EUR/MW over the 12-month test period, compared to 85,500 for the na\"{i}ve and essentially zero for the long-only reference. The trading performance is robust to cost assumptions (cost drag below 2\% across scenarios) and consistent across all calendar months.

\section{Recommendations}

Based on our findings, we offer the following recommendations for practitioners:

\begin{itemize}
  \item \textbf{Use Random Forest as the primary forecasting model} for the French day-ahead market. It outperforms both the na\"{i}ve benchmark and XGBoost on risk-adjusted metrics, and its MDI feature importances provide interpretable insights into price drivers.

  \item \textbf{Prioritise fundamental features in stable regimes; retain weather capability for crisis regimes.} Under normal market conditions, model development effort is better spent improving the quality of load forecasts, fuel price integration, and cross-border flow data rather than adding more meteorological variables. However, the 2022 result demonstrates that weather features provide statistically significant value during crisis regimes characterised by low nuclear availability and decoupled gas-weather co-movement. A production model should include weather features with a regime-detection mechanism that activates or de-weights them based on observable market conditions.

  \item \textbf{Implement periodic retraining.} The current static model is trained through April 2024. A monthly or quarterly retraining schedule would adapt to structural market changes (increasing renewable penetration, nuclear fleet evolution) and is a prerequisite for production deployment.

  \item \textbf{Complement point forecasts with prediction intervals.} The high RMSE relative to MAE indicates the presence of large forecast errors during spike periods. Risk managers should complement model predictions with scenario analysis for extreme price events.
\end{itemize}

\section{Broader Implications}

The regime-dependent weather finding reflects a broader insight about
information aggregation in commodity markets. In stable conditions, when market
prices (TTF gas) and operational forecasts (RTE load forecast) already
incorporate weather expectations, raw meteorological variables become
statistically redundant. This principle — that feature selection should account
for the information already embedded in price-based variables — is likely
applicable beyond the electricity market to other commodity forecasting
problems. However, the 2022 result shows that this redundancy is conditional:
structural supply shocks that break the normal co-movement between commodity
prices and weather can reinstate the direct weather channel. Forecasters should
monitor the stability of these co-movement relationships and treat weather
feature inclusion as a dynamic decision rather than a permanent design choice.

From a market design perspective, this thesis illustrates how the quality of
publicly available data (ENTSO-E generation data, RTE load forecasts) can
substitute for proprietary weather modelling infrastructure during normal market
conditions, lowering the barriers to participation in day-ahead electricity
price forecasting. During crisis regimes, however, direct weather integration
provides measurable value that cannot be obtained from public market data alone.

\section{Code and Data Availability}

The complete codebase for this thesis — including data collection pipelines (ENTSO-E API, ERA5 CDS API, fuel prices), feature engineering, model training, evaluation, and the trading backtest — is available upon request. Data from ENTSO-E and ERA5 are publicly available under their respective open data licences.
